\documentclass[11pt,a4paper,oneside]{article}

% -------------------------------------------------
% Basic packages
% -------------------------------------------------
\usepackage[T1]{fontenc}
\usepackage[utf8]{inputenc}
\usepackage[spanish,es-noquoting,es-nodecimaldot]{babel}
\usepackage{lmodern}
\usepackage{microtype}
\usepackage[a4paper,margin=1.18in,headheight=15pt]{geometry}
\usepackage{amsmath,amssymb,amsfonts,amsthm,mathtools}
\usepackage{bm}
\usepackage{enumitem}
\usepackage{booktabs}
\usepackage{array}
\usepackage{xcolor}
\usepackage{hyperref}
\usepackage{setspace}
\usepackage{fancyhdr}
\usepackage{titlesec}
\usepackage[most]{tcolorbox}
\usepackage[capitalize,noabbrev]{cleveref}

\hypersetup{
  colorlinks=true,
  linkcolor=blue!50!black,
  citecolor=red!55!black,
  urlcolor=teal!60!black,
  pdftitle={SIAMDA: Sistema Integrado de Análisis Multirresolución y Desempeño Académico},
  pdfauthor={César Galindo López y Gabriel Villafuerte C.}
}

\theoremstyle{definition}
\newtheorem{definition}{Definición}[section]
\newtheorem{remark}[definition]{Observación}

\definecolor{ihesblue}{RGB}{25,62,114}
\definecolor{ihesgold}{RGB}{160,120,40}
\pagestyle{fancy}
\fancyhf{}
\fancyhead[L]{\textsc{SIAMDA}}
\fancyhead[R]{\textsc{Galindo López \& Villafuerte C.}}
\fancyfoot[C]{\thepage}
\renewcommand{\headrulewidth}{0.35pt}
\renewcommand{\footrulewidth}{0pt}
\titleformat{\section}{\Large\bfseries\color{ihesblue}}{\thesection.}{0.65em}{}
\titleformat{\subsection}{\large\bfseries\color{ihesblue}}{\thesubsection.}{0.65em}{}
\titlespacing*{\section}{0pt}{1.9em}{1.05em}
\titlespacing*{\subsection}{0pt}{1.45em}{0.95em}
\tcbset{colback=white,colframe=ihesblue!65!black,boxrule=0.5pt,arc=1.5pt}
\setlength{\parskip}{0.85em}
\setstretch{1.13}
\setlength{\parindent}{0pt}

\newtcolorbox{ideaBox}[1][]{enhanced,breakable,colback=blue!2,colframe=ihesblue,
  borderline west={1.5pt}{0pt}{ihesblue},title=Idea central,fonttitle=\bfseries,#1}
\newtcolorbox{summaryBox}[1][]{enhanced,breakable,colback=gray!4,colframe=black!55,
  borderline west={1.5pt}{0pt}{black!70},title=Fórmulas clave,fonttitle=\bfseries,#1}
\newtcolorbox{warningBox}[1][]{enhanced,breakable,colback=red!2,colframe=red!55!black,
  borderline west={1.5pt}{0pt}{red!75!black},title=Alcance y limitaciones,fonttitle=\bfseries,#1}
\newtcolorbox{implBox}[1][]{enhanced,breakable,
  colback=violet!3!white,colframe=violet!50!black,
  borderline west={2pt}{0pt}{violet!65!black},
  title=Implementación,fonttitle=\bfseries\color{violet!80!black},
  attach boxed title to top left={yshift=-2mm,xshift=4mm},
  boxed title style={colback=violet!12,colframe=violet!50!black,boxrule=0.5pt},#1}

\title{\vspace{-1.5cm}\textbf{\Large SIAMDA}\\[2pt]\large Sistema Integrado de Análisis Multirresolución y Desempeño Académico}
\author{César Galindo López \\ \small Universidad Panamericana \textperiodcentered\ Facultad de Empresariales, Ciudad UP
  \and Gabriel Villafuerte C. \\ \small \href{https://github.com/Gabriel44svg}{github.com/Gabriel44svg}}
\date{\small 2026}

\begin{document}
\maketitle

\begin{ideaBox}
SIAMDA combina un pipeline cuantitativo (calificaciones, parciales, estado de riesgo académico)
con un pipeline cualitativo (respuestas abiertas de encuesta) para producir un índice de
sentimiento por alumno mediante \emph{embeddings} semánticos y la transformada discreta de
onduleta (DWT). La eficiencia estadística de ese índice, entendido como estimador de un
parámetro poblacional, se certifica comparando su varianza empírica contra la cota de
Cramér--Rao bajo dos modelos de referencia (gaussiano y Laplaciano). Esta nota documenta la
arquitectura del sistema y la formulación matemática exacta implementada en el código
(\texttt{processing/datos.py}, \texttt{processing/nlp\_wavelet.py},
\texttt{processing/cramer\_rao.py}), sin reportar corridas experimentales sobre datos reales de
curso: es una descripción de método, no un estudio de caso.
\end{ideaBox}

\section{Motivación}

Un curso genera dos tipos de señal sobre el desempeño de un alumno: una \emph{cuantitativa}
(calificaciones parciales, tareas, exposición) y una \emph{cualitativa} (lo que el alumno
escribe en una encuesta de retroalimentación abierta). La señal cuantitativa es la que
normalmente se usa para detectar riesgo académico; la cualitativa suele leerse de forma
impresionista, sin métrica. SIAMDA propone una métrica para la segunda señal —un índice de
sentimiento $\hat\theta\in[-1,1]$ por alumno, obtenido de sus respuestas de texto— y pregunta
si esa métrica aporta información adicional sobre el riesgo académico, vía su correlación
empírica con la calificación final.

\section{Arquitectura del sistema}

El sistema se organiza en cuatro módulos:

\begin{enumerate}[leftmargin=1.4em]
  \item \textbf{Ingesta y normalización} (\texttt{datos.py}): lectura de archivos de
  calificaciones (CSV/Excel) con detección automática de columnas (número de cuenta, nombre,
  parciales \texttt{E1,E2,\dots}, examen ponderado, tareas, exposición, validación), manejo de
  texto mixto en celdas no numéricas (\emph{``te presentas a final''}, \emph{``está en
  reposición''}, \emph{``baja''}, etc.) y cómputo de una bandera binaria de riesgo.
  \item \textbf{NLP y extracción wavelet} (\texttt{nlp\_wavelet.py}): generación de
  \emph{embeddings} de oración multilingües sobre las respuestas abiertas de la encuesta, y
  descomposición multirresolución de cada vector de embedding vía DWT.
  \item \textbf{Validación teórica} (\texttt{cramer\_rao.py}): cómputo de la información de
  Fisher y la cota de Cramér--Rao para el índice de sentimiento agregado, bajo los modelos
  gaussiano y Laplaciano, más una prueba de normalidad (Shapiro--Wilk) y una estimación
  bootstrap de la varianza empírica del estimador.
  \item \textbf{Visualización interactiva}: interfaz Streamlit con páginas para carga de datos,
  métricas del curso, análisis de sentimiento y validación teórica, con interpretaciones
  automáticas en lenguaje natural para cada métrica.
\end{enumerate}

\section{Bandera de riesgo académico}

Sea $c_i$ la calificación final del alumno $i$ y $e_i\in\{\text{Normal},\ldots\}$ su estado de
validación (reposición, examen final, baja, etc.). El sistema define
\[
r_i = \mathbf{1}\{c_i < 70\} \ \vee\ \mathbf{1}\{e_i \neq \text{Normal}\},
\]
es decir, un alumno está en riesgo si su calificación cae por debajo del umbral aprobatorio
(70) o si tiene un estado especial registrado, independientemente de la calificación numérica.

\section{Índice de sentimiento por transformada wavelet}

\subsection{Embeddings y descomposición}

Para cada alumno $i$, sea $x_i\in\mathbb{R}^d$ el embedding semántico de su respuesta de texto,
obtenido con un modelo de \emph{sentence-transformers} multilingüe (\texttt{paraphrase-
multilingual-MiniLM-L12-v2}), normalizado a norma unitaria. Tratando $x_i$ como una señal
discreta en $\mathbb{R}^d$, se aplica la transformada discreta de onduleta (DWT) con familia
Daubechies~4 y nivel de descomposición $J$:
\[
\mathrm{wavedec}(x_i, J) = \big(c^{(i)}_A,\ c^{(i)}_{D,1},\ \ldots,\ c^{(i)}_{D,J}\big),
\]
donde $c^{(i)}_A$ son los coeficientes de aproximación (contenido semántico global, de baja
frecuencia) y $c^{(i)}_{D,j}$ los coeficientes de detalle en el nivel $j$ (variación local, de
alta frecuencia). Se definen las energías
\[
E^{\mathrm{apr}}_i = \big\lVert c^{(i)}_A \big\rVert_2^2,
\qquad
E^{\mathrm{det}}_i = \sum_{j=1}^{J} \big\lVert c^{(i)}_{D,j} \big\rVert_2^2 .
\]

\begin{remark}
La energía de aproximación captura contenido contextual estable (ideas organizadas,
vocabulario consistente); la energía de detalle actúa como proxy de variación local abrupta en
el embedding, interpretada aquí como carga emocional aguda en el texto.
\end{remark}

\subsection{Índice de sentimiento}

El índice de sentimiento del alumno $i$ se define como
\[
\hat\theta_i \;=\; \tanh\!\left(\frac{E^{\mathrm{det}}_i - E^{\mathrm{apr}}_i}
{E^{\mathrm{det}}_i + E^{\mathrm{apr}}_i + \varepsilon}\right), \qquad \varepsilon = 10^{-9},
\]
de modo que $\hat\theta_i\in(-1,1)$: valores cercanos a $-1$ indican energía concentrada en la
aproximación (respuesta contextual, poca carga emocional local) y valores cercanos a $+1$
indican energía concentrada en los detalles (alta variabilidad emocional localizada).

\section{Validación estadística: información de Fisher y cota de Cramér--Rao}

Sea $\hat\theta = (\hat\theta_1,\ldots,\hat\theta_N)$ la muestra de índices de sentimiento del
grupo, tratada como $N$ observaciones i.i.d.\ de un parámetro poblacional $\theta$ (el
sentimiento promedio del grupo), y $\bar\theta = \frac1N\sum_i \hat\theta_i$ su estimador de
media muestral. La cota de Cramér--Rao acota inferiormente la varianza de cualquier estimador
insesgado de $\theta$:
\[
\mathrm{Var}(\bar\theta) \;\ge\; \frac{1}{F(\theta)},
\qquad
F(\theta) = -\,\mathbb{E}\!\left[\frac{\partial^2}{\partial\theta^2}\ln p(X;\theta)\right].
\]

\begin{summaryBox}
\textbf{Modelo gaussiano.} Bajo $\theta\sim\mathcal N(\mu,\sigma^2)$, con
$\hat\mu=\bar\theta$ y $\hat\sigma^2$ la varianza muestral insesgada:
\[
F(\mu) = \frac{N}{\hat\sigma^2}, \qquad
\mathrm{CRB}_\mu = \frac{1}{F(\mu)} = \frac{\hat\sigma^2}{N}.
\]
\textbf{Modelo Laplaciano.} Bajo $\theta\sim\mathrm{Laplace}(\mu,b)$, con $\hat\mu=\mathrm{mediana}(\hat\theta)$
y $\hat b = \frac1N\sum_i |\hat\theta_i-\hat\mu|$:
\[
F(\mu) = \frac{N}{\hat b^2}, \qquad
\mathrm{CRB}_\mu = \frac{\hat b^2}{N}.
\]
\textbf{Eficiencia empírica del estimador (EES).} Con $\widehat{\mathrm{Var}}(\bar\theta)$
estimada por \emph{bootstrap} ($B=500$ remuestreos con reemplazo):
\[
\mathrm{EES} = \min\!\left(\frac{\mathrm{CRB}_\mu}{\widehat{\mathrm{Var}}(\bar\theta)},\ 1\right).
\]
\end{summaryBox}

El sistema complementa este cómputo con una prueba de normalidad de Shapiro--Wilk sobre
$\hat\theta$: si el valor $p$ resultante es menor a $0.05$, se recomienda reportar el modelo
Laplaciano (más robusto a colas pesadas) en lugar del gaussiano, y usar la mediana muestral
—el estimador óptimo bajo ese modelo— en vez de la media.

\begin{remark}
La cota de Cramér--Rao certifica algo preciso y limitado: qué tan cerca está la \emph{media
muestral} de $\hat\theta$ de ser un estimador óptimo de sí misma bajo un modelo paramétrico
supuesto. No certifica que $\hat\theta$ mida ``sentimiento real'' correctamente, ni sustituye
una validación externa (p.\,ej.\ contra anotación humana) del índice.
\end{remark}

\section{Correlación con desempeño académico}

Como paso final del análisis, el sistema calcula el coeficiente de correlación de Pearson $r$
entre $\hat\theta_i$ y la calificación $c_i$ de cada alumno, junto con su valor $p$ asociado y
el tamaño de muestra $N$. Una correlación negativa significativa ($r<0$, $p<0.05$) se interpreta
como evidencia de que mayor carga emocional en las respuestas se asocia con menor desempeño;
una correlación positiva significativa se interpreta como posible motivación extrínseca. El
sistema clasifica la fuerza de la asociación en umbrales ($|r|\ge 0.6$ fuerte, $|r|\ge 0.3$
moderada, $|r|<0.3$ débil) y explícitamente reporta cuando la correlación no es significativa,
sin forzar una narrativa de riesgo donde no hay evidencia estadística.

\begin{warningBox}
SIAMDA es una herramienta de apoyo a la decisión docente, no un instrumento diagnóstico
automatizado. El índice $\hat\theta$ depende del modelo de embeddings elegido, de la familia y
nivel de la wavelet, y del tamaño de la encuesta; con $N<30$ las estimaciones bootstrap de
varianza son inestables y el propio sistema advierte de ello en la interfaz. Las
interpretaciones textuales que genera el módulo \texttt{interpretacion.py} son sugerencias de
lectura, no conclusiones definitivas, y deben cotejarse con el criterio del docente y, cuando
el caso lo amerite, con apoyo psicopedagógico institucional.
\end{warningBox}

\begin{implBox}
Interfaz Streamlit (\texttt{app.py}) con cuatro páginas modulares
(\texttt{pages/carga.py}, \texttt{pages/metricas.py}, \texttt{pages/sentimiento.py},
\texttt{pages/validacion.py}); \texttt{sentence-transformers} para embeddings,
\texttt{PyWavelets} para la DWT, \texttt{scipy.stats} para Shapiro--Wilk y correlación de
Pearson, \texttt{plotly} para visualización interactiva (incluyendo histograma bootstrap de
$\bar\theta$ con bandas $\pm\sqrt{\mathrm{CRB}}$, comparación de varianzas, y gráfica Q--Q de
$\hat\theta$ contra la normal). Código y datos de ejemplo en
\href{https://github.com/Gabriel44svg/Galindo-C.-J.-y-Villafuerte-C.-G.-2026-.-SIAMDA}{github.com/Gabriel44svg/Galindo-C.-J.-y-Villafuerte-C.-G.-2026-.-SIAMDA}.
\end{implBox}

\section{Conclusión}

SIAMDA integra dos señales de un curso —calificaciones y respuestas abiertas de encuesta— bajo
un mismo tablero, y añade a la segunda una métrica reproducible ($\hat\theta$, vía embeddings y
DWT) junto con su certificado estadístico de eficiencia (información de Fisher y cota de
Cramér--Rao, bajo modelos gaussiano y Laplaciano). El resultado es un sistema de alerta
temprana que hace explícito, en vez de impresionista, el vínculo entre el tono de las
respuestas de un grupo y su desempeño, y que reporta honestamente cuándo esa correlación no es
estadísticamente significativa.

\end{document}
